\frfilename{mt136.tex}
\versiondate{22.6.05}
\copyrightdate{2000}

\def\chaptername{Pelengkap}
\def\sectionname{Teorema Kelas Monoton}

\newsection{*136}

Pada bagian terakhir jilid ini, saya menyajikan dua teorema mengenai
aljabar-$\sigma$ beserta beberapa akibat sederhana. Hasil-hasil ini ditempatkan
di sini karena saya tidak menemukan tempat alami untuknya dalam Jilid 2.
Walaupun keduanya (terutama 136B) merupakan bagian dari teknik dasar teori
ukuran dan memiliki banyak penerapan yang tersebar luas, keduanya tidak sentral
bagi pendekatan khusus yang saya pilih dan, jika Anda menghendakinya, dapat
dikesampingkan sampai diperlukan.

\leader{136A}{Lema} Misalkan $X$ suatu himpunan, dan $\Cal A$ suatu keluarga
subhimpunan $X$. Maka pernyataan-pernyataan berikut ekuivalen:

\inset{(i) $X\in\Cal A$, $B\setminus A\in\Cal A$ kapan pun $A$,
$B\in\Cal A$ dan $A\subseteq B$, dan $\bigcup_{n\in\Bbb N}A_n\in\Cal A$
kapan pun $\sequencen{A_n}$ merupakan barisan tak-menurun dalam $\Cal A$;

(ii) $\emptyset\in\Cal A$, $X\setminus A\in\Cal A$ untuk setiap
$A\in\Cal A$, dan $\bigcup_{n\in\Bbb N}A_n\in\Cal A$ kapan pun
$\sequencen{A_n}$ merupakan barisan saling lepas dalam $\Cal A$.}

\proof{{\bf (i)$\Rightarrow$(ii)} Andaikan (i) benar. Tentu saja
$\emptyset=X\setminus X$ termasuk dalam $\Cal A$ dan
$X\setminus A\in\Cal A$ untuk setiap $A\in\Cal A$.
Jika $A$, $B\in\Cal A$ saling lepas, maka
$A\subseteq X\setminus B\in\Cal A$, sehingga
$(X\setminus B)\setminus A$ dan komplemennya $A\cup B$ termasuk dalam
$\Cal A$. Jadi, jika $\sequencen{A_n}$ merupakan barisan saling lepas dalam
$\Cal A$, $\bigcup_{i\le n}A_i\in\Cal A$ untuk setiap $n$, dan
$\bigcup_{n\in\Bbb N}A_n$ adalah gabungan suatu barisan tak-menurun dalam
$\Cal A$, sehingga termasuk dalam $\Cal A$. Dengan demikian (ii) benar.

\medskip

{\bf (ii)$\Rightarrow$(i)} Jika (ii) benar, maka tentu saja
$X=X\setminus\emptyset$ termasuk dalam $\Cal A$. Jika $A$ dan $B$ adalah
anggota $\Cal A$ sedemikian sehingga $A\subseteq B$, maka $X\setminus B$
termasuk dalam $\Cal A$ dan saling lepas dengan $A$, sehingga
$A\cup(X\setminus B)$ dan komplemennya $B\setminus A$ termasuk dalam
$\Cal A$. Dengan demikian klausa kedua (i) terpenuhi. Untuk klausa ketiga,
jika $\sequencen{A_n}$ merupakan barisan tak-menurun dalam $\Cal A$, maka
$A_0,A_1\setminus A_0,A_2\setminus A_1,\ldots$ merupakan barisan saling lepas
dalam $\Cal A$, sehingga gabungannya $\bigcup_{n\in\Bbb N}A_n$ termasuk
dalam $\Cal A$.
}%end of proof of 136A

\medskip

\noindent{\bf Definisi} Jika $\Cal A\subseteq\Cal PX$ memenuhi syarat-syarat
(i) dan/atau (ii) di atas, keluarga tersebut disebut suatu {\bf kelas Dynkin} dari
subhimpunan-subhimpunan $X$.

\leader{136B}{Teorema Kelas Monoton} Misalkan $X$ suatu himpunan dan
$\Cal A$ suatu kelas Dynkin dari subhimpunan-subhimpunan $X$. Andaikan bahwa
$\Cal I\subseteq\Cal A$ sedemikian sehingga $I\cap J\in\Cal I$ untuk semua
$I$, $J\in\Cal I$. Maka $\Cal A$ memuat aljabar-$\sigma$ dari subhimpunan-
subhimpunan $X$ yang dibangkitkan oleh $\Cal I$.

\proof{{\bf (a)} Misalkan $\frak S$ keluarga kelas-kelas Dynkin dari
subhimpunan-subhimpunan $X$ yang memuat $\Cal I$. Maka mudah diperiksa,
dengan menggunakan (i) atau (ii) dalam 136A, bahwa irisan
$\Sigma=\bigcap\frak S$ juga merupakan suatu kelas Dynkin (bandingkan 111Ga).
Karena $\Cal A\in\frak S$, $\Sigma\subseteq\Cal A$.

\medskip

{\bf (b)} Jika $H\in\Sigma$, maka

\Centerline{$\Sigma_H=\{E:E\in\Sigma,\,E\cap H\in\Sigma\}$}

\noindent merupakan suatu kelas Dynkin. \Prf\ ($\alpha$)
$X\cap H=H\in\Sigma$ sehingga $X\in\Sigma_H$. ($\beta$) Jika $A$,
$B\in\Sigma_H$ dan $A\subseteq B$, maka $A\cap H$, $B\cap H$ termasuk
dalam $\Sigma$ dan $A\cap H\subseteq B\cap H$; akibatnya

\Centerline{$(B\setminus A)\cap H=(B\cap H)\setminus(A\cap H)\in\Sigma$}

\noindent dan $B\setminus A\in\Sigma_H$. ($\gamma$) Jika
$\sequencen{A_n}$ merupakan barisan tak-menurun dalam $\Sigma_H$, maka
$\sequencen{A_n\cap H}$ merupakan barisan tak-menurun dalam $\Sigma$, sehingga

\Centerline{$(\bigcup_{n\in\Bbb N}A_n)\cap H
=\bigcup_{n\in\Bbb N}(A_n\cap H)\in\Sigma$}

\noindent dan $\bigcup_{n\in\Bbb N}A_n\in\Sigma_H$.\ \Qed

Akibatnya, jika $I\cap H\in\Sigma$ untuk setiap $I\in\Cal I$, sehingga
$\Sigma_H\supseteq\Cal I$, maka $\Sigma_H\in\frak S$ dan harus sama dengan
$\Sigma$.

\medskip

{\bf (c)} Berikutnya kita dapati bahwa $G\cap H\in\Sigma$ untuk semua $G$,
$H\in\Sigma$. \Prf\ Ambil $I$, $J\in\Cal I$. Kita tahu bahwa
$I\cap J\in\Cal I$. Karena $I$ sembarang, $\Sigma_J=\Sigma$ dan
$H\in\Sigma_J$, yakni $H\cap J\in\Sigma$. Karena $J$ sembarang,
$\Sigma_H=\Sigma$ dan $G\in\Sigma_H$, yakni $G\cap H\in\Sigma$.\ \Qed

\medskip

{\bf (d)} Karena $\Sigma$ merupakan kelas Dynkin,
$\emptyset=X\setminus X\in\Sigma$. Selain itu,

\Centerline{$G\cup H
=X\setminus((X\setminus G)\cap(X\setminus H))\in\Sigma$}

\noindent untuk sembarang $G$, $H\in\Sigma$ (menggunakan (c)).
Jadi, jika $\sequencen{G_n}$ merupakan sembarang barisan dalam $\Sigma$,
$G'_n=\bigcup_{i\le n}G_i\in\Sigma$ untuk setiap $n$ (dengan induksi pada
$n$). Namun, $\sequencen{G'_n}$ sekarang merupakan barisan tak-menurun dalam
$\Sigma$, sehingga

\Centerline{$\bigcup_{n\in\Bbb N}G_n
=\bigcup_{n\in\Bbb N}G'_n\in\Sigma$.}

Ini berarti bahwa $\Sigma$ memenuhi semua syarat dalam 111A dan merupakan
suatu aljabar-$\sigma$ dari subhimpunan-subhimpunan $X$. Karena
$\Cal I\subseteq\Sigma$, $\Sigma$ harus memuat aljabar-$\sigma$ $\Sigma'$ dari
subhimpunan-subhimpunan $X$ yang dibangkitkan oleh $\Cal I$. Jadi
$\Sigma'\subseteq\Sigma\subseteq\Cal A$, sebagaimana diperlukan.

(Sebenarnya, tentu saja, $\Sigma=\Sigma'$, karena $\Sigma'\in\frak S$.)
}%end of proof of 136B

\cmmnt{\medskip

\noindent{\bf Catatan} Saya pernah melihat hasil ini disebut
{\bf Teorema Kelas Sierpi\'nski} dan
{\bf Teorema $\pmb{\pi}$-$\pmb{\lambda}$}.}

\leader{136C}{Akibat} Misalkan $X$ suatu himpunan, dan $\mu$, $\nu$ dua
ukuran yang didefinisikan pada $X$ dengan domain masing-masing $\Sigma$,
$\Tau$. Andaikan bahwa $\mu X=\nu X<\infty$, dan bahwa
$\Cal I\subseteq\Sigma\cap\Tau$ merupakan keluarga himpunan sedemikian
sehingga $\mu I=\nu I$ untuk setiap $I\in\Cal I$ dan
$I\cap J\in\Cal I$ untuk semua $I$, $J\in\Cal I$. Maka $\mu E=\nu E$ untuk
setiap $E$ dalam aljabar-$\sigma$ dari subhimpunan-subhimpunan $X$ yang
dibangkitkan oleh $\Cal I$.

\proof{ Intinya adalah bahwa

\Centerline{$\Cal A=\{H:H\in\Sigma\cap\Tau,\,\mu H=\nu H\}$}

\noindent merupakan suatu kelas Dynkin dari subhimpunan-subhimpunan $X$.
\Prf\ Saya menggunakan (ii) dalam 136A. Tentu saja
$\emptyset\in\Cal A$. Jika $A\in\Cal A$, maka

\Centerline{$\mu(X\setminus A)=\mu X-\mu A=\nu X-\nu A
=\nu(X\setminus A)$}

\noindent (karena $\mu X=\nu X<\infty$, sehingga pengurangannya aman), dan
$X\setminus A\in\Cal A$. Jika $\sequencen{A_n}$ merupakan barisan saling lepas
dalam $\Cal A$ dan $A=\bigcup_{n\in\Bbb N}A_n$, maka

\Centerline{$\mu A=\sum_{n=0}^{\infty}\mu A_n=\sum_{n=0}^{\infty}\nu A_n
=\nu A$,}

\noindent dan $\bigcup_{n\in\Bbb N}A_n\in\Cal A$.\ \Qed

Karena $\Cal I\subseteq\Cal A$, 136B memberi tahu kita bahwa
aljabar-$\sigma$ $\Sigma'$ yang dibangkitkan oleh $\Cal I$ termuat dalam
$\Cal A$, yakni, $\mu$ dan $\nu$ sama pada $\Sigma'$.
}%end of proof of 136C

\vleader{72pt}{136D}{Akibat} Misalkan $\mu$, $\nu$ dua ukuran pada
$\BbbR^r$,
dengan $r\ge 1$, yang keduanya didefinisikan, dan bernilai sama, pada semua
interval berbentuk

\Centerline{$\ocint{-\infty,a}=\{x:x\le a\}=\{(\xi_1,\ldots,\xi_r):
\xi_i\le\alpha_i$ untuk setiap $i\le r\}$}

\noindent untuk $a=(\alpha_1,\ldots,\alpha_r)\in\BbbR^r$. Andaikan pula bahwa
$\mu\BbbR^r<\infty$. Maka $\mu$ dan $\nu$ bernilai sama pada semua
subhimpunan Borel dari $\BbbR^r$.

\proof{ Dalam 136C, ambil $X=\BbbR^r$ dan $\Cal I$ sebagai himpunan
interval-interval $\ocint{-\infty,a}$. Maka $I\cap J\in\Cal I$ untuk semua
$I$, $J\in\Cal I$, karena
$\ocint{-\infty,a}\cap\ocint{-\infty,b}=\ocint{-\infty,a\wedge b}$, dengan
menuliskan
$a\wedge b=(\min(\alpha_1,\beta_1),\ldots,\min(\alpha_r,\beta_r))$ jika
$a=(\alpha_1,\ldots,\alpha_r)$, $b=(\beta_1,\ldots,\beta_r)\in\BbbR^r$.
Selain itu, dengan menetapkan $\tbf{n}=(n,\ldots,n)$ untuk $n\in\Bbb N$,

\Centerline{$\nu\BbbR^r=\lim_{n\to\infty}\nu\ocint{-\infty,\tbf{n}}
=\lim_{n\to\infty}\mu\ocint{-\infty,\tbf{n}}=\mu\BbbR^r$.}

\noindent Jadi semua syarat 136C terpenuhi dan $\mu$, $\nu$ bernilai sama
pada aljabar-$\sigma$ $\Sigma$ yang dibangkitkan oleh $\Cal I$. Namun ini
tepat merupakan aljabar-$\sigma$ himpunan-himpunan Borel, menurut 121J.
}%end of proof of 136D

\leader{136E}{Aljabar himpunan: Definisi} Misalkan $X$ suatu himpunan.
Suatu keluarga $\Cal E\subseteq\Cal PX$ merupakan suatu {\bf aljabar} atau
{\bf medan} dari subhimpunan-subhimpunan $X$ jika

\inset{(i) $\emptyset\in\Cal E$;

(ii) untuk setiap $E\in\Cal E$, komplemennya $X\setminus E$ termasuk
dalam $\Cal E$;

(iii) untuk setiap $E$, $F\in\Cal E$, $E\cup F\in\Cal E$.}

\leader{136F}{Catatan-catatan}\cmmnt{ {\bf (a)} Saya dapat saja
memperkenalkan gagasan ini dalam Bab 11, bersama `aljabar-$\sigma$'. Saya tidak
memperkenalkannya di sana, kecuali dalam beberapa latihan, karena tampaknya
sudah cukup banyak definisi baru dalam \S111, dan karena tidak ada hal
substansial yang ingin saya katakan mengenai aljabar himpunan.

\header{136Fb}}{\bf (b)} Jika $\Cal E$ suatu aljabar dari subhimpunan-
subhimpunan $X$, maka

\Centerline{$E\cap F\cmmnt{\mskip5mu =X\setminus((X\setminus
E)\cup(X\setminus
F))}$,
\quad$E\setminus F\cmmnt{\mskip5mu =E\cap(X\setminus F)}$,}

\Centerline{$E_0\cup E_1\cup\ldots\cup E_n$,
\quad$E_0\cap E_1\cap\ldots\cap E_n$}

\noindent termasuk dalam $\Cal E$ untuk semua $E$, $F$,
$E_0,\ldots,E_n\in\Cal E$. \cmmnt{(Gunakan induksi pada $n$ untuk yang
terakhir.)}

\spheader 136Fc Suatu aljabar-$\sigma$ dari subhimpunan-subhimpunan $X$
\cmmnt{(tentu saja)} merupakan suatu aljabar dari subhimpunan-subhimpunan $X$.

\vleader{42pt}{136G}{Teorema} Misalkan $X$ suatu himpunan dan $\Cal E$ suatu
aljabar dari subhimpunan-subhimpunan $X$. Andaikan bahwa
$\Cal A\subseteq\Cal PX$ merupakan keluarga himpunan sedemikian sehingga

\inset{($\alpha$) $\bigcup_{n\in\Bbb N}A_n\in\Cal A$ untuk setiap barisan
tak-menurun $\sequencen{A_n}$ dalam $\Cal A$,

($\beta$) $\bigcap_{n\in\Bbb N}A_n\in\Cal A$ untuk setiap barisan
tak-menaik $\sequencen{A_n}$ dalam $\Cal A$,

($\gamma$) $\Cal E\subseteq\Cal A$.}

\noindent Maka $\Cal A$ memuat aljabar-$\sigma$ dari subhimpunan-subhimpunan
$X$ yang dibangkitkan oleh $\Cal E$.

\proof{ Saya menggunakan gagasan yang sama seperti dalam 136B.

\medskip

{\bf (a)} Misalkan $\frak S$ keluarga semua himpunan
$\Cal S\subseteq\Cal PX$ yang memenuhi ($\alpha$)-($\gamma$). Maka irisannya
$\Sigma=\bigcap\frak S$ juga memenuhi syarat-syarat tersebut. Karena
$\Cal A\in\frak S$, $\Sigma\subseteq\Cal A$.

\medskip

{\bf (b)} Jika $H\in\Sigma$, maka

\Centerline{$\Sigma_H=\{E:E\in\Sigma,\,E\cap H\in\Sigma\}$}

\noindent memenuhi syarat-syarat ($\alpha$)-($\beta$). \Prf\ ($\alpha$)
Jika $\sequencen{A_n}$ merupakan barisan tak-menurun dalam $\Sigma_H$, maka
$\sequencen{A_n\cap H}$ merupakan barisan tak-menurun dalam $\Sigma$, sehingga

\Centerline{$(\bigcup_{n\in\Bbb N}A_n)\cap H
=\bigcup_{n\in\Bbb N}(A_n\cap
H)\in\Sigma$}

\noindent dan $\bigcup_{n\in\Bbb N}A_n\in\Sigma_H$. ($\beta$)
Demikian pula, jika $\sequencen{A_n}$ merupakan barisan tak-menaik dalam
$\Sigma_H$, maka $\bigcap_{n\in\Bbb N}A_n\cap H\in\Sigma$ sehingga
$\bigcap_{n\in\Bbb N}A_n\in\Sigma_H$.\ \Qed

Akibatnya, jika $E\cap H\in\Sigma$ untuk setiap $E\in\Cal E$, sehingga
$\Sigma_H$ juga memenuhi ($\gamma$), maka $\Sigma_H\in\frak S$ dan harus
sama dengan $\Sigma$.
\medskip

{\bf (c)} Akibatnya $G\cap H\in\Sigma$ untuk semua $G$, $H\in\Sigma$.
\Prf\ Ambil $E$, $F\in\Cal E$. Kita tahu bahwa $E\cap F\in\Cal E$. Karena
$E$ sembarang, $\Sigma_F=\Sigma$ dan $H\in\Sigma_F$, yakni
$H\cap F\in\Sigma$. Karena $F$ sembarang, $\Sigma_H=\Sigma$ dan
$G\in\Sigma_H$, yakni $G\cap H\in\Sigma$.\ \Qed

\medskip

{\bf (d)} Selanjutnya, $\Sigma^*=\{X\setminus H:H\in\Sigma\}\in\frak S$.
\Prf\ ($\alpha$) Jika $\sequencen{A_n}$ merupakan barisan tak-menurun dalam
$\Sigma^*$, maka $\sequencen{X\setminus A_n}$ merupakan barisan tak-menaik
dalam $\Sigma$, sehingga

\Centerline{$\bigcup_{n\in\Bbb N}A_n=X
\setminus\bigcap_{n\in\Bbb N}(X\setminus A_n)\in\Sigma^*$.}

\noindent ($\beta$) Demikian pula, jika $\sequencen{A_n}$ merupakan barisan
tak-menaik dalam $\Sigma^*$, maka

\Centerline{$\bigcap_{n\in\Bbb N}A_n
=X\setminus\bigcup_{n\in\Bbb N}(X\setminus A_n)\in\Sigma^*$.}

\noindent ($\gamma$) Jika $E\in\Cal E$, maka $X\setminus E\in\Cal E$
sehingga $X\setminus E\in\Sigma$ dan $E\in\Sigma^*$.\ \QeD\ Akibatnya
$\Sigma\subseteq\Sigma^*$, yakni $X\setminus H\in\Sigma$ untuk setiap
$H\in\Sigma$.

\medskip

{\bf (e)} Dengan menggabungkan (c) dan (d) dengan fakta bahwa
$X\in\Sigma$ (karena $X\in\Cal E$) dan bahwa gabungan suatu barisan
tak-menurun dalam $\Sigma$ termasuk dalam $\Sigma$ (menurut syarat
($\alpha$)), kita melihat bahwa argumen yang sama seperti dalam bagian (d)
bukti 136B menunjukkan bahwa $\Sigma$ merupakan suatu aljabar-$\sigma$ dari
subhimpunan-subhimpunan $X$. Jadi, seperti dalam 136B, kita menyimpulkan bahwa
aljabar-$\sigma$ yang dibangkitkan oleh $\Cal E$ termuat dalam $\Sigma$ dan
karena itu dalam $\Cal A$.
}%end of proof of 136G

\leader{*136H}{Proposisi} Misalkan $(X,\Sigma,\mu)$ suatu ruang ukur
sedemikian sehingga $\mu X<\infty$, dan $\Cal E$ suatu subaljabar dari
$\Sigma$; misalkan $\Sigma'$ merupakan aljabar-$\sigma$ dari subhimpunan-
subhimpunan $X$ yang dibangkitkan oleh $\Cal E$. Jika $F\in\Sigma'$ dan
$\epsilon>0$, terdapat $E\in\Cal E$ sedemikian sehingga
$\mu(F\symmdiff E)\le\epsilon$.

\proof{ Misalkan $\Cal A$ keluarga himpunan $F\in\Sigma$ sedemikian sehingga

\Centerline{untuk setiap $\epsilon>0$ terdapat $E\in\Cal E$ sedemikian
sehingga $\mu(F\symmdiff E)\le\epsilon$.}

\noindent Maka $\Cal A$ merupakan suatu kelas Dynkin. \Prf\ Saya memeriksa
ketiga syarat dalam 136A(i). ($\alpha$) $X\in\Cal A$ karena
$X\in\Cal E$. ($\beta$) Jika $F_1$, $F_2\in\Cal A$ dan $\epsilon>0$,
terdapat $E_1$, $E_2\in\Cal E$ sedemikian sehingga
$\mu(F_i\symmdiff E_i)\le\bover12\epsilon$ untuk kedua $i$; sekarang
$E_1\setminus E_2\in\Cal E$ dan

\Centerline{$(F_1\setminus F_2)\symmdiff(E_1\setminus E_2)
\subseteq(F_1\symmdiff E_1)\cup(F_2\symmdiff E_2)$,}

\noindent sehingga

\Centerline{$\mu((F_1\setminus F_2)\symmdiff(E_1\setminus E_2))
\le\mu(F_1\symmdiff E_1)+\mu(F_2\symmdiff E_2)\le\epsilon$.}

\noindent Karena $\epsilon$ sembarang, $F_1\setminus F_2\in\Cal A$.
($\gamma$) Jika $\sequencen{F_n}$ merupakan barisan tak-menurun dalam
$\Cal A$, dengan gabungan $F$, dan $\epsilon>0$, maka

\Centerline{$\lim_{n\to\infty}\mu F_n=\mu F\le\mu X<\infty$,}

\noindent sehingga terdapat $n\in\Bbb N$ sedemikian sehingga
$\mu(F\setminus F_n)\le\bover12\epsilon$. Sekarang terdapat
$E\in\Cal E$ sedemikian sehingga
$\mu(F_n\symmdiff E)\le\bover12\epsilon$; karena
$F\symmdiff E\subseteq(F\setminus F_n)\cup(F_n\symmdiff E)$,
$\mu(F\symmdiff E)\le\epsilon$. Karena $\epsilon$ sembarang,
$F\in\Cal A$.\ \Qed

Karena $\Cal E\subseteq\Cal A$ dan $\Cal E$ tertutup terhadap $\cap$,
$\Cal A$ memuat aljabar-$\sigma$ $\Sigma'$ yang dibangkitkan oleh $\Cal E$,
sebagaimana diklaim.
}%end of proof of 136H

\exercises{
\leader{136X}{Latihan dasar $\pmb{>}$(a)}
%\spheader 136Xa
Misalkan $X$ suatu himpunan dan $\Cal A$ suatu keluarga subhimpunan $X$.
Tunjukkan bahwa pernyataan-pernyataan berikut ekuivalen:

\quad(i) $X\in\Cal A$ dan $B\setminus A\in\Cal A$ kapan pun $A$,
$B\in\Cal A$ dan $A\subseteq B$;

\quad(ii) $\emptyset\in\Cal A$, $X\setminus A\in\Cal A$ untuk setiap
$A\in\Cal A$ dan $A\cup B\in\Cal A$ kapan pun $A$,
$B\in\Cal A$ saling lepas.
%136A

\spheader 136Xb Andaikan bahwa $X$ suatu himpunan dan
$\Cal A\subseteq\Cal PX$. Tunjukkan bahwa $\Cal A$ merupakan suatu
aljabar-$\sigma$ dari subhimpunan-subhimpunan $X$ jika dan hanya jika
keluarga itu merupakan suatu kelas Dynkin dan $A\cap B\in\Cal A$ kapan pun $A$,
$B\in\Cal A$.
%136B

\spheader 136Xc\dvArevised{2010}
Misalkan $X$ suatu himpunan, dan $\Cal I$ suatu keluarga subhimpunan
$X$ sedemikian sehingga $I\cap J\in\Cal I$ untuk semua $I$, $J\in\Cal I$;
misalkan $\Sigma$ merupakan aljabar-$\sigma$ dari subhimpunan-subhimpunan
$X$ yang dibangkitkan oleh $\Cal I$.
Misalkan pula $\mu$ dan $\nu$ dua ukuran pada $X$, keduanya berdomain
$\Sigma$, sedemikian sehingga $\mu I=\nu I<\infty$ untuk setiap
$I\in\Cal I$.
Tunjukkan bahwa
$\mu E=\nu E$ kapan pun $E\in\Sigma$ ditutupi oleh suatu barisan dalam
$\Cal I$.
({\it Petunjuk\/}: Untuk $J\in\Cal I$, tetapkan
$\mu_JE=\mu(E\cap J)$,
$\nu_JE=\nu(E\cap J)$ untuk $E\in\Sigma$. Gunakan 136C untuk menunjukkan
bahwa $\mu_J=\nu_J$ untuk setiap $J$.)
%136C

\sqheader 136Xd Tetapkan $X=\{0,1,2,3\}$,
$\Cal I=\{X,\{0,1\},\{0,2\}\}$. Temukan dua ukuran berbeda $\mu$, $\nu$
pada $X$, keduanya didefinisikan pada aljabar-$\sigma$ $\Cal PX$ dan dengan
$\mu I=\nu I<\infty$ untuk setiap $I\in\Cal I$.
%136C

\spheader 136Xe Misalkan $\Sigma$ keluarga subhimpunan
$\coint{0,1}$ yang dapat dinyatakan sebagai gabungan berhingga interval-
interval setengah terbuka $\coint{a,b}$. Tunjukkan bahwa $\Sigma$ merupakan
suatu aljabar dari subhimpunan-subhimpunan $\coint{0,1}$.
%136E

\spheader 136Xf Misalkan $X$ suatu himpunan, dan $\Cal I$ suatu keluarga
subhimpunan $X$ sedemikian sehingga $I\cap J\in\Cal I$ kapan pun $I$,
$J\in\Cal I$. Misalkan $\Sigma$ keluarga himpunan terkecil sedemikian sehingga
$X\in\Sigma$, $F\setminus E\in\Sigma$ kapan pun $E$, $F\in\Sigma$ dan
$E\subseteq F$, dan $\Cal I\subseteq\Sigma$. Tunjukkan bahwa $\Sigma$
merupakan suatu aljabar dari subhimpunan-subhimpunan $X$.
%136E

\spheader 136Xg Misalkan $X$ suatu himpunan, dan
$\Cal E$ suatu aljabar dari subhimpunan-subhimpunan $X$. Suatu fungsional
$\nu:\Cal E\to\Bbb R$ disebut {\bf aditif} ({\bf berhingga}) jika
$\nu(E\cup F)=\nu E+\nu F$ kapan pun $E$,
$F\in\Cal E$ dan $E\cap F=\emptyset$. (i) Tunjukkan bahwa dalam kasus ini
$\nu(E\cup F)+\nu(E\cap F)=\nu E+\nu F$ untuk semua $E$, $F\in\Cal E$.
(ii) Tunjukkan bahwa jika $\nu E\ge 0$ untuk setiap $E\in\Cal E$, maka
$\nu(\bigcup_{i\le n}E_i)\le\sum_{i=0}^n\nu E_i$ untuk semua
$E_0,\ldots,E_n\in\Cal E$.
%136F

\sqheader 136Xh Misalkan $X$ suatu himpunan, dan $\Cal A$ suatu keluarga
subhimpunan $X$ sedemikian sehingga ($\alpha$) $\emptyset$, $X$ termasuk
dalam $\Cal A$ ($\beta$) $A\cap B\in\Cal A$ untuk semua $A$,
$B\in\Cal A$ ($\gamma$) $A\cup B\in\Cal A$ kapan pun $A$, $B\in\Cal A$ dan
$A\cap B=\emptyset$. Tunjukkan bahwa
$\{A:A\in\Cal A,\,X\setminus A\in\Cal A\}$ merupakan suatu aljabar dari
subhimpunan-subhimpunan $X$.
%136F

\sqheader 136Xi Misalkan $X$ suatu himpunan, dan $\Cal A$ suatu keluarga
subhimpunan $X$ sedemikian sehingga ($\alpha$) $\emptyset$, $X$ termasuk
dalam $\Cal A$ ($\beta$) $\bigcap_{n\in\Bbb N}A_n\in\Cal A$ untuk setiap
barisan $\sequencen{A_n}$ dalam $\Cal A$ ($\gamma$)
$\bigcup_{n\in\Bbb N}A_n\in\Cal A$ untuk setiap barisan saling lepas
$\sequencen{A_n}$ dalam $\Cal A$. Tunjukkan bahwa
$\{A:A\in\Cal A,\,X\setminus A\in\Cal A\}$ merupakan suatu aljabar-$\sigma$
dari subhimpunan-subhimpunan $X$.
%136Xh, 136F

\sqheader 136Xj Misalkan $\Cal A$ suatu keluarga subhimpunan $\Bbb R$
sedemikian sehingga (i) $\bigcap_{n\in\Bbb N}A_n\in\Cal A$ untuk setiap
barisan $\sequencen{A_n}$ dalam $\Cal A$ (ii)
$\bigcup_{n\in\Bbb N}A_n\in\Cal A$ untuk setiap barisan saling lepas
$\sequencen{A_n}$ dalam $\Cal A$ (iii) setiap interval terbuka
$\ooint{a,b}$ termasuk dalam $\Cal A$. Tunjukkan bahwa setiap subhimpunan
Borel dari $\Bbb R$ termasuk dalam $\Cal A$. ({\it Petunjuk\/}: tunjukkan
bahwa setiap interval setengah terbuka $\coint{a,b}$, $\ocint{a,b}$ termasuk
dalam $\Cal A$, dan karena itu semua interval $\ocint{-\infty,a}$,
$\coint{a,\infty}$; sekarang gunakan 136Xi.)
%136Xi, 136F

\sqheader 136Xk Misalkan $X$ suatu himpunan, $\Cal E$ suatu aljabar dari
subhimpunan-subhimpunan $X$, dan $\Cal A$ suatu keluarga subhimpunan $X$
sedemikian sehingga ($\alpha$) $\bigcap_{n\in\Bbb N}A_n\in\Cal A$ untuk
setiap barisan tak-menaik $\sequencen{A_n}$ dalam $\Cal A$ ($\beta$)
$\bigcup_{n\in\Bbb N}A_n\in\Cal A$ untuk setiap barisan saling lepas dalam
$\Cal A$ ($\gamma$) $\Cal E\subseteq\Cal A$. Tunjukkan bahwa
aljabar-$\sigma$ himpunan yang dibangkitkan oleh $\Cal E$ termuat dalam
$\Cal A$. \Hint{gunakan metode 136B untuk mereduksi ke kasus ketika
$A\cap B\in\Cal A$ untuk setiap $A$, $B\in\Cal A$; sekarang gunakan 136Xi.}
%136Xi, 136F

\leader{136Y}{Latihan lanjutan (a)} Misalkan $X$ suatu himpunan dan
$\Cal E$ suatu aljabar dari subhimpunan-subhimpunan $X$. Misalkan
$\nu:\Cal E\to\coint{0,\infty}$ suatu fungsional nonnegatif yang aditif dalam
pengertian 136Xg. Definisikan $\theta:\Cal PX\to\coint{0,\infty}$ dengan
menetapkan

\Centerline{$\theta A=\inf\{\sum_{n=0}^{\infty}\nu E_n:\sequencen{E_n}$
merupakan barisan dalam $\Cal E$ yang menutupi $A\}$}

\noindent untuk setiap $A\subseteq X$. (i) Tunjukkan bahwa $\theta$ merupakan
suatu ukuran luar pada $X$ dan bahwa $\theta E\le\nu E$ untuk setiap
$E\in\Cal E$. (ii) Misalkan $\mu$ ukuran pada $X$ yang didefinisikan dari
$\theta$ dengan metode \Caratheodory, dan $\Sigma$ domainnya. Tunjukkan bahwa
$\Cal E\subseteq\Sigma$ dan bahwa $\mu E\le\nu E$ untuk setiap
$E\in\Cal E$. (iii) Tunjukkan bahwa pernyataan-pernyataan berikut ekuivalen:
($\alpha$) $\mu E=\nu E$
untuk setiap $E\in\Cal E$ ($\beta$) $\theta X=\nu X$ ($\gamma$) kapan pun
$\sequencen{E_n}$ merupakan barisan tak-menaik dalam $\Cal E$ dengan irisan
kosong, $\lim_{n\to\infty}\nu E_n=0$.

\spheader 136Yb Misalkan $X$ suatu himpunan, $\Cal E$ suatu aljabar dari
subhimpunan-subhimpunan $X$, dan $\nu$ suatu fungsional aditif nonnegatif pada
$\Cal E$. Misalkan $\Sigma$ merupakan aljabar-$\sigma$ dari subhimpunan-
subhimpunan $X$ yang dibangkitkan oleh $\Cal E$. Tunjukkan bahwa terdapat
paling banyak satu ukuran $\mu$ pada $X$ dengan domain $\Sigma$ yang
memperluas $\nu$, dan bahwa ukuran demikian ada jika dan hanya jika
$\lim_{n\to\infty}\nu E_n=0$ untuk setiap barisan tak-menaik
$\sequencen{E_n}$ dalam $\Cal E$ dengan irisan kosong.

\spheader 136Yc Misalkan $X$ suatu himpunan. Misalkan $\Cal G$ suatu keluarga
subhimpunan $X$ sedemikian sehingga (i) $G\cap H\in\Cal G$ untuk semua $G$,
$H\in\Cal G$ (ii) untuk setiap $G\in\Cal G$ terdapat suatu barisan
$\sequencen{G_n}$ dalam $\Cal G$ sedemikian sehingga
$X\setminus G=\bigcap_{n\in\Bbb N}G_n$. Misalkan $\Cal A$ suatu keluarga
subhimpunan $X$ sedemikian sehingga ($\alpha$) $\emptyset$, $X\in\Cal A$
($\beta$) $\bigcap_{n\in\Bbb N}A_n\in\Cal A$
untuk setiap barisan tak-menaik $\sequencen{A_n}$ dalam
$\Cal A$ ($\gamma$) $\bigcup_{n\in\Bbb N}A_n\in\Cal A$ untuk setiap
barisan saling lepas dalam $\Cal A$ ($\delta$) $\Cal G\subseteq\Cal A$.
Tunjukkan bahwa aljabar-$\sigma$ himpunan yang dibangkitkan oleh $\Cal G$
termuat dalam $\Cal A$.
%136Xk, 136F
}%end of exercises

\endnotes{
\Notesheader{136} Hasil paling berguna di sini adalah 136B; hasil ini akan
diperlukan dalam Bab 27, dan berguna pada berbagai bagian lain dalam Jilid 2,
sering kali melalui akibat-akibatnya 136C dan 136Xc. Tentu saja 136C, seperti
akibatnya 136D dan kasus khususnya 136Yb, hanya dapat digunakan secara langsung
pada ukuran yang tidak mengambil nilai $\infty$, karena kita harus mengetahui
bahwa $\mu(F\setminus E)=\mu F-\mu E$ untuk himpunan terukur
$E\subseteq F$; itulah sebabnya hasil ini menjadi menonjol hanya ketika kita
mengkhususkan diri pada ukuran probabilitas (yang seluruh ruangnya berukuran
$1$). Karena itu saya menyertakan 136Xc untuk menunjukkan teknik yang dapat
membawa kita selangkah lebih jauh. Saya rasa kita belum benar-benar siap untuk
ukuran umum pada himpunan-himpunan Borel dari $\BbbR^r$, tetapi saya
menyebutkan 136D untuk menunjukkan jenis kelas $\Cal I$ yang dapat muncul
dalam 136B.

Kedua teorema di sini (136B, 136G) menjawab pertanyaan: jika diberikan suatu
keluarga himpunan $\Cal I$, operasi apa yang harus kita lakukan untuk membangun
aljabar-$\sigma$ $\Sigma$ yang dibangkitkan oleh $\Cal I$? Untuk
$\Cal I$ sembarang, tentu saja kita mengharapkan perlunya komplemen dan gabungan
barisan. Inti teorema-teorema di sini adalah bahwa jika $\Cal I$ mempunyai
sejumlah struktur tertentu, kita dapat mencapai $\Sigma$ dengan operasi yang
lebih terbatas; jadi, jika $\Cal I$ merupakan suatu aljabar himpunan, maka
gabungan dan irisan {\it monoton} sudah cukup (136G). Tentu saja terdapat tak
terhitung banyaknya variasi atas tema ini. Saya menawarkan 136Xh–136Xj sebagai
hasil khas yang benar-benar akan digunakan dalam Jilid 4, dan 136Xk serta 136Yc
sebagai contoh modifikasi yang mungkin. Terdapat versi abstrak 136B dalam 313G
di Jilid 3.

Setelah mulai mempertimbangkan perluasan suatu aljabar himpunan menjadi suatu
aljabar-$\sigma$, wajar untuk menanyakan syarat-syarat yang memungkinkan suatu
fungsional pada aljabar himpunan diperluas menjadi suatu ukuran. Syarat
aditivitas (136Xg) jelas diperlukan, dan hampir sama jelasnya bahwa syarat itu
tidak mencukupi.
Saya menyertakan 136Ya–136Yb sebagai syarat perlu dan cukup yang paling penting
di antara banyak syarat agar suatu fungsional aditif dapat diperluas menjadi
suatu ukuran. Kita harus kembali ke hal ini dalam Jilid 4.
}%end of comment

\frnewpage
